\documentclass[../../main.tex]{subfiles}
\begin{document}

\begin{center}
  {\fontsize{48pt}{58pt}\selectfont \textbf{1971}\normalsize\textbf{年度}}
\end{center}

\begin{table}[htb]
  \centering
  \caption{出題テーマと難易度}
  \begin{tabular}{|c|p{9cm}|c|}
    \hline
    \textbf{問題番号} & \textbf{テーマ} & \textbf{難易度} \\
    \hline
    第1問 & 格子点と極限 & \\
    \hline
    第2問 & 積分不等式の最小定数 & \\
    \hline
    第3問 & 積分と三角関数の最大値 & \\
    \hline
    第4問 & 座標変換と面積 & \\
    \hline
    第5問 & 不等式証明と極限 & \\
    \hline
    第6問 & 確率と漸化式 & \\
    \hline
  \end{tabular}
\end{table}

\begin{table}[htb]
  \centering
  \caption{解答}
  \renewcommand{\arraystretch}{1.4}
  \begin{tabular}{|c|c|p{9cm}|}
    \hline
    \textbf{問題番号} & \textbf{小問} & \textbf{答え} \\
    \hline
    第1問 & - & 証明問題 \\
    \hline
    第2問 & - & $K=\dfrac{1}{2}a$ \\
    \hline
    第3問 & - & $\theta=0,\pi,2\pi$ \\
    \hline
    第4問 & (1) & 図示，$S(\theta)=\dfrac{1}{6}s^3+\dfrac{1}{2}sc^2-\dfrac{1}{2}cs^2$（$0<\theta\le\pi/4$，$c,s$入替で他方） \\
    \cline{2-3}
    & (2) & 証明問題 \\
    \hline
    第5問 & (1) & 証明問題 \\
    \cline{2-3}
    & (2) & $+0$ \\
    \hline
    第6問 & (1) & $p_{n+1}=(\alpha+\beta-1)p_n+(1-\beta)$ \\
    \cline{2-3}
    & (2) & $\dfrac{1-\beta}{2-(\alpha+\beta)}$ \\
    \hline
  \end{tabular}
\end{table}

\paragraph{試験形態}
6問．試験時間は要確認．

\end{document}
